% ============================================================================
% KONFIGURASI TEMPLATE — ubah nilai di file ini saja.
% ============================================================================

% Pilihan jenjang yang valid: sarjana atau magister
\def\jenjang{sarjana}
\def\jenjangSarjana{sarjana}
\def\jenjangMagister{magister}
\newif\ifmagister
\ifx\jenjang\jenjangMagister
  \magistertrue
\else\ifx\jenjang\jenjangSarjana
  \magisterfalse
\else
  \PackageError{UII-Template}{Jenjang '\jenjang' tidak dikenali}{Gunakan nilai 'sarjana' atau 'magister' pada config.tex.}
\fi\fi

% Data karya ilmiah
\newcommand{\judulTA}{BAGIAN INI ADALAH BAGIAN JUDUL}
\newcommand{\subjudulTA}{TULIS SUBJUDUL DI SINI JIKA ADA}
\newcommand{\judulInggris}{ENGLISH THESIS TITLE}
\newcommand{\subjudulInggris}{ENGLISH SUBTITLE IF APPLICABLE}

% Data mahasiswa
\newcommand{\namaMahasiswa}{NAMA MAHASISWA}
\newcommand{\nimMahasiswa}{20523000}
\newcommand{\konsentrasiMagister}{NAMA KONSENTRASI}

% Data pembimbing. Kosongkan namaPembimbingDua jika hanya ada satu pembimbing.
\newcommand{\namaPembimbing}{Nama Dosen Pembimbing, S.T., M.Eng.}
\newcommand{\nipPembimbing}{NIK/NIP Pembimbing}
\newcommand{\namaPembimbingDua}{}
\newcommand{\nipPembimbingDua}{}

% Data penguji
\newcommand{\namaPengujiSatu}{Nama Dosen Penguji 1, S.T., M.Sc.}
\newcommand{\nipPengujiSatu}{NIK/NIP Penguji 1}
\newcommand{\namaPengujiDua}{Nama Dosen Penguji 2, S.T., M.Cs.}
\newcommand{\nipPengujiDua}{NIK/NIP Penguji 2}
\newcommand{\namaPengujiTiga}{Nama Dosen Penguji 3, S.T., M.Kom.}
\newcommand{\nipPengujiTiga}{NIK/NIP Penguji 3}

% Data program studi dan waktu
\newcommand{\ketuaProgramStudi}{Nama Ketua Program Studi, S.T., M.Sc.}
\newcommand{\nipKetuaProgramStudi}{NIK/NIP Ketua Program Studi}
\newcommand{\tanggalPengesahan}{Tanggal Bulan Tahun}
\newcommand{\bulanLulus}{Bulan}
\newcommand{\tahunLulus}{2026}
\newcommand{\kotaTerbit}{Yogyakarta}

% Kata kunci
\newcommand{\kataKunci}{kata kunci pertama, kata kunci kedua, kata kunci ketiga}
\newcommand{\keywords}{keyword one, keyword two, keyword three}

% Nilai turunan otomatis berdasarkan jenjang.
\ifmagister
  \newcommand{\jenisDokumen}{Tesis}
  \newcommand{\jenisDokumenKapital}{TESIS}
  \newcommand{\namaProgram}{PROGRAM MAGISTER}
  \newcommand{\namaProgramKalimat}{Program Studi Informatika Program Magister}
  \newcommand{\gelarTarget}{Magister Komputer}
  \newcommand{\singkatanGelar}{M.Kom.}
\else
  \newcommand{\jenisDokumen}{Tugas Akhir}
  \newcommand{\jenisDokumenKapital}{TUGAS AKHIR}
  \newcommand{\namaProgram}{PROGRAM SARJANA}
  \newcommand{\namaProgramKalimat}{Program Studi Informatika Program Sarjana}
  \newcommand{\gelarTarget}{Sarjana Komputer}
  \newcommand{\singkatanGelar}{S.Kom.}
\fi
